\newtoggle{withbonus}
\settoggle{withbonus}{false} % set true to show bonus points in grade table

% setting underlying course name
\newcommand{\mycourse}{Informatik}

% name of the class
\newcommand{\myclass}{\faPeace\faIceCream~2. Klassen~\faIceCream\faPeace}

% set date
\newcommand{\mydate}{Freitag, 26. Juni 2026}

\title{\textbf{randomisierte Algorithmen, Graphen, Datenintegrität}}
\author{Thomas Graf\\
\mycourse, \myclass}
\date{\mydate}
\maketitle
\thispagestyle{headandfoot}

\begin{center}
	\fbox{\parbox{140mm}{\centering
			\begin{itemize}
				\item Dauer der Prüfung: $45$ Minuten
				\item erlaubte Hilfsmittel:
				\begin{itemize}
					\item Schreibzeug
					\item \textbf{ohne Taschenrechner}
				\end{itemize}
			\end{itemize}}}
\end{center}

\vspace{10mm}

\begin{center}
	Vorname: \rule{45mm}{0.8pt}\qquad Nachname: \rule{45mm}{0.8pt}
\end{center}

\vspace{20mm}

\begin{center}
	\iftoggle{withbonus}{
		\combinedgradetable[h][questions]
	}{
		\gradetable[h][questions]
	}
\end{center}

\vspace{15mm}

\begin{center}
	Note:
	\vspace{10mm}

	\rule{8mm}{1.2pt}
\end{center}

\clearpage


\begin{questions}
\question
\begin{parts}
  \part[1] Es sei $x := 0011011$ ein binärer String. Bestimmen Sie $\text{Nummer}(x)$. 
  \part[1] Bestimmen Sie $\text{Bin}(25)$.
  \part[1] Für einen binären String $x$ gilt 
  \begin{align*}
   \text{Nummer}(x) = 2^k - 1.
  \end{align*}
  Dabei ist $k$ eine positive natürliche Zahl. Bestimmen Sie $\abs{x}$ in Abhängigkeit von $k$.
  \part[1] Bestimmen Sie $\text{Prim}(40)$.
    \part[1] Ronald Rivest möchte unser randomisierte Kommunikationsprotokoll für die beiden Datensätze $x := 001110$ und $y := 1110$ durchführen. Was ist hier das Problem? Welchen Rat würden Sie Donald geben?
\end{parts}

\question[2] In Zürich sei der Datensatz $x := 001111$ auf dem Rechner $R_1$ abgespeichert. $R_1$ wird das Kommunikationsprotokoll mit einem Rechner $R_2$ in Boston durchführen.
Dabei sei $y$ der Datensatz, welcher auf $R_2$ abgespeichert ist. In Zürich wurde zufällig die Primzahl $p := 7$ gewählt. Geben Sie ein konkretes Beispiel eines (sinnvollen) Datensatzes $y$, sodass das Protokoll falsch entscheiden wird.

\question Gegeben seien zwei Datensätze $x$ und $y$ mit einer jeweiligen Länge von $10^{24}$ Bits. 
\begin{parts}
\part[2] Geben Sie eine vernünftige obere Schranke für den Kommunikationsaufwand unseres Kommunikationsprotokolls zum Abgleich von $x$ und $y$ an.
\part[1] Wie viele Bits müsste ein deterministischer Algorithmus zur Beantwortung der Frage $x\stackrel{?}{=}y$ mindestens kommunizieren?
\end{parts}
\droptotalpoints

\question[1]
Gegeben sind zwei Funktionen $f$ und $g$, welche auf ganz $\N$ definiert sind und $g(n) \neq 0$ für alle $n\in\N$. Es gelte die Gleichung
\begin{align}\label{eq:limit}
  \lim_{n\to\infty}\frac{f(n)}{g(n)} = 1.
\end{align}
Ist es korrekt, dass \cref{eq:limit} die Gleichung
\begin{align*}
  \lim_{n\to\infty} \abs{f(n) - g(n)} = 0
\end{align*}
impliziert? Begründen Sie die Behauptung oder finden Sie ein Gegenbeispiel.

\clearpage

\question
\textbf{Graphentheorie: Eulerwege}

Gegeben ist ein ungerichteter Graph mit den Knoten
\[
	V = \{A, B, C, D, E\}
\]
und den Kanten
\[
	E = \{\{A,B\}, \{A,C\}, \{B,C\}, \{B,D\}, \{B,E\}, \{C,D\}, \{C,E\}, \{D,E\}\}.
\]

\begin{parts}
	\part[1] Bestimmen Sie die Grade der Knoten $A$ und $B$. \vspace{20mm}
	\part[1] Existiert in diesem Graphen ein Eulerweg? Begründen Sie Ihre Antwort kurz mit Hilfe der Knotengrade. \vspace{20mm}
\end{parts}
\droptotalpoints

\begin{solution}
\begin{enumerate}
	\item Grad($A$) = 2 (verbunden mit $B, C$). \\
	Grad($B$) = 4 (verbunden mit $A, C, D, E$).
	\item Ein Eulerweg existiert genau dann, wenn es höchstens zwei Knoten mit ungeradem Grad gibt.
	Prüfen wir die restlichen Grade:
	Grad($C$) = 4, aber Grad($D$) = 3 und Grad($E$) = 3.
	Da es genau zwei Knoten ($D$ und $E$) mit ungeradem Grad gibt, existiert \textbf{ein} Eulerweg.
\end{enumerate}
\end{solution}

\clearpage

\question[3]
\textbf{Dijkstra-Algorithmus vervollständigen}
Vervollständigen Sie den in \cref{listing:Dijkstra} gezeigten Pseudocode des Dijkstra-Algorithmus, indem Sie die fehlenden Teile einfügen. Die fehlenden Teile sind mit Punkten markiert. Die Variablen $d_s[u]$ und $c(u,v)$ bezeichnen die Distanz von Knoten $s$ zu Knoten $u$ und die Kosten der Kante von Knoten $u$ zu Knoten $v$.
\begin{lstlisting}[language=Python,numbers=none,escapeinside={(*}{*)},caption=\texttt{Dijkstra-Algorithmus},label=listing:Dijkstra]
### Dijkstra-Algorithmus ###

## Initialisierung ##
für alle Knoten u des Graphen setze:
    (*$d_s[u] = .....$*)
(*$d_s[s] \leftarrow 0$*)
(*$S \leftarrow \emptyset$*)
(*$Q \leftarrow $*) Menge aller Knoten des Graphen

## Schritte ##
solange ..... nicht leer ist, mach:
    (*$u\leftarrow$*) Knoten aus (*$Q$*) mit kleinstem (*$d_s[u]$*)
    (*$Q \leftarrow Q \setminus \lrc{u}$*)
    (*$S\leftarrow S \cup \lrc{u}$*)

    für alle von (*$u$*) direkt erreichbaren Knoten (*$v$*) mach:
        # ist der Weg von (*$s$*) nach (*$v$*) kürzer über (*$u$*)?
        falls ......................... dann:
            (*$d_s[v] \leftarrow d_s[u] + c(u,v)$*)
\end{lstlisting}


\clearpage


\question
\textbf{Fehlererkennung und -korrektur}

Gegeben ist eine binäre Kodierung mit den folgenden drei gültigen Codewörtern:
\begin{itemize}
	\item $C_1 = 00000$
	\item $C_2 = 00111$
	\item $C_3 = 11100$
\end{itemize}

\begin{parts}
	\part[1] Bestimmen Sie den minimalen Abstand dieser Kodierung. \vspace{15mm}
	\part[1] Wie viele Bitfehler können mit dieser Kodierung \textit{sicher erkannt} werden? \vspace{15mm}
	\part[1] Wie viele Bitfehler können mit dieser Kodierung \textit{sicher korrigiert} werden? \vspace{15mm}
\end{parts}

\begin{solution}
\begin{enumerate}
	\item Die Abstände sind
	\begin{align*}
		d(C_1, C_2) &= 3 \\
		d(C_1, C_3) &= 3 \\
		d(C_2, C_3) &= 4
	\end{align*}
	Der minimale Abstand ist also \(d_{\min} = 3\).

	\item Sicher erkannt werden können
	\[
		d_{\min} - 1 = 2
	\]
	Bitfehler.

	\item Sicher korrigiert werden können
	\[
		\left\lfloor \frac{d_{\min} - 1}{2} \right\rfloor = \left\lfloor \frac{2}{2} \right\rfloor = 1
	\]
	Bitfehler.
\end{enumerate}
\end{solution}

\end{questions}